\section{Results}
\label{sec:results}

This section reports the outcomes of our three research questions.
All results are drawn from \texttt{results/summary.csv} and \texttt{results/full\_analysis.ipynb}, which were produced by \texttt{evaluation/integrated\_statistics.py} before any data were viewed.

% -----------------------------------------------------------------------
\subsection{Executability}
\label{subsec:exec_results}

Of the 25~Java generated test suites, 24~(96\%) were classified as \texttt{EXECUTABLE}.
The single non-executable suite, \texttt{JAVA\_016\_parseCsvLine}, failed at compile time because the generated code referenced the \texttt{unit\_id} string as a class name rather than \texttt{JavaAlgorithms}.
This suite is retained with \texttt{NA} for coverage and mutation in all primary analyses.

Of the 25~Python generated test suites, 22~(88\%) were executable.
Three suites (\texttt{PY\_013}, \texttt{PY\_020}, \texttt{PY\_025}) contained syntax errors and are classified as \texttt{TEST\_RUNTIME\_ERROR}.

% -----------------------------------------------------------------------
\subsection{RQ1 --- Branch Coverage vs.\ 75\% Target}
\label{subsec:rq1}

Table~\ref{tab:rq1} summarises the branch coverage outcomes and sign-test results.

\begin{table}[ht]
\centering
\caption{RQ1: Branch Coverage Results (Primary Complete-Case Analysis)}
\label{tab:rq1}
\begin{tabular}{lcccccc}
\toprule
Lang & $n$ valid & Median & IQR & $k$ above 75\% & $p_{\text{raw}}$ & $p_{\text{Holm}}$ \\
\midrule
Java   & 24/25 & 95.45\% & [89.62, 100.00] & 23/24 & $1.49 \times 10^{-6}$ & $2.98 \times 10^{-6}$ \\
Python & 22/25 & 100.00\% & [87.85, 100.00] & 19/22 & $4.28 \times 10^{-4}$ & $4.28 \times 10^{-4}$ \\
\bottomrule
\end{tabular}
\end{table}

For Java, the exact right-tailed sign test yields $k = 23$ out of 24~non-tied observations above the 75\% threshold ($p_{\text{raw}} = 1.49 \times 10^{-6}$; Holm-adjusted $p = 2.98 \times 10^{-6}$).
For Python, $k = 19$ out of 22~non-tied observations exceed the threshold ($p_{\text{raw}} = 4.28 \times 10^{-4}$; Holm-adjusted $p = 4.28 \times 10^{-4}$).
Both adjusted $p$-values are below $\alpha = 0.05$; we \textbf{reject $H_0$} for both languages.

\textbf{Worst-case sensitivity.}
When the non-executable suites are recoded as $0\%$ branch coverage, the Java sign test yields $p = 9.72 \times 10^{-6}$ and the Python sign test yields $p = 7.32 \times 10^{-3}$; both remain below $\alpha = 0.05$, confirming the result is robust to non-executability.

% -----------------------------------------------------------------------
\subsection{RQ2 --- GPT-4o vs.\ Randoop Mutation Score (Java)}
\label{subsec:rq2}

Table~\ref{tab:rq2} reports the paired mutation-score comparison.
Of the 25~Java units, 24~provide paired data; \texttt{JAVA\_016\_parseCsvLine} contributes no GPT mutation score and is excluded from the pairing.

\begin{table}[ht]
\centering
\caption{RQ2: Paired Mutation Score --- GPT-4o vs.\ Randoop (Java, $n = 24$)}
\label{tab:rq2}
\begin{tabular}{lcc}
\toprule
Statistic & GPT-4o & Randoop \\
\midrule
Median mutation score     & 86.61\%        & 15.38\%       \\
IQR                       & [80.00, 95.21] & [6.25, 63.16] \\
Median difference ($d_i$) & \multicolumn{2}{c}{$+62.14$ pp} \\
IQR of difference         & \multicolumn{2}{c}{[17.08, 86.16] pp} \\
Bootstrap 95\% CI         & \multicolumn{2}{c}{[31.58, 84.62] pp} \\
Wilcoxon $W$              & \multicolumn{2}{c}{190.0} \\
$p$-value (one-sided)     & \multicolumn{2}{c}{$6.58 \times 10^{-5}$} \\
\bottomrule
\end{tabular}
\end{table}

The one-sided paired Wilcoxon signed-rank test yields $W = 190.0$, $p = 6.58 \times 10^{-5} < 0.05$; we \textbf{reject $H_0$}.
The median difference of $+62.14$~percentage points with a 95\% CI of $[31.58, 84.62]$~pp confirms that the advantage is both statistically significant and practically large.
The sign-test sensitivity analysis (19~positive pairs, 0~negative, 5~ties; $p = 1.91 \times 10^{-6}$) is consistent with the primary result.

% -----------------------------------------------------------------------
\subsection{RQ3 --- Cyclomatic Complexity vs.\ Test Quality}
\label{subsec:rq3}

Table~\ref{tab:rq3} reports the Spearman correlation results.
Python mutation score is excluded from the confirmatory family because pytest-mutagen produced no mutants for any Python function; the three remaining correlations are Holm-corrected jointly.

\begin{table}[ht]
\centering
\caption{RQ3: Spearman Correlation --- CC vs.\ Test Quality (Holm-corrected, 3 testable pairs)}
\label{tab:rq3}
\begin{tabular}{llcccc}
\toprule
Language & Outcome & $n$ & $\rho$ & $p_{\text{raw}}$ & $p_{\text{Holm}}$ \\
\midrule
Java   & Branch Coverage & 24 & $-0.517$ & $0.0049$ & $0.0146$ \\
Java   & Mutation Score  & 24 & $+0.109$ & $0.695$  & $0.695$  \\
Python & Branch Coverage & 22 & $-0.512$ & $0.0074$ & $0.0148$ \\
Python & Mutation Score  & 0  & N/A      & N/A      & EXCL.    \\
\bottomrule
\end{tabular}
\end{table}

For Java branch coverage, Spearman $\rho = -0.517$ ($p_{\text{Holm}} = 0.0146 < 0.05$); we \textbf{reject $H_0$}.
For Python branch coverage, Spearman $\rho = -0.512$ ($p_{\text{Holm}} = 0.0148 < 0.05$); we \textbf{reject $H_0$}.
For Java mutation score, $\rho = +0.109$ ($p_{\text{Holm}} = 0.695$); we \textbf{fail to reject $H_0$}.
Python mutation score is \textbf{not testable} due to the tool limitation described in Section~\ref{subsec:metrics}.
